%%%%%%%%%%%%%%%%%%%%%%%%%%%%%%%%%%%%%%%%%%%%%%%%%%%%%%%%%%%%%%%%%%%%%%%%%%%%%%%%%%
\begin{frame}[fragile]\frametitle{}
\begin{center}
{\Large आरोग्य व दीर्घायुष्य --- सी. शिवशंकरन}
\end{center}

{\tiny (Ref: सी. शिवशंकरन (Aircel चे माजी संस्थापक) --- आरोग्य व दीर्घायुष्य, YouTube)}

\end{frame}

%%%%%%%%%%%%%%%%%%%%%%%%%%%%%%%%%%%%%%%%%%%%%%%%%%%%%%%%%%%
\begin{frame}[fragile]\frametitle{रोगमुक्त दीर्घायुष्य शक्य आहे}

	\begin{itemize}
	\item ९५ वर्षापर्यंत मधुमेह, कर्करोग, हृदयविकाराशिवाय जगणे शक्य आहे
	\item येथे `रोगमुक्त' म्हणजे खोकला-सर्दी नव्हे --- गंभीर आजार नसणे
	\item हे दैनंदिन निवडींवर अवलंबून --- काय करतो, काय खातो, कसे राहतो
	\item आनुवंशिकता बंदूक भरते --- जीवनशैली ट्रिगर दाबते
	\item जीवनशैली सुमारे ९३\% परिणामांवर नियंत्रण ठेवते
	\end{itemize}

{\tiny (Ref: सी. शिवशंकरन --- दीर्घायुष्य)}

\end{frame}

%%%%%%%%%%%%%%%%%%%%%%%%%%%%%%%%%%%%%%%%%%%%%%%%%%%%%%%%%%%
\begin{frame}[fragile]\frametitle{पवित्र त्रिकूट: खा, श्वास घे, झोप}

	\begin{itemize}
	\item हळू, लांब श्वास सोडणे = नर्व्हस सिस्टमला शांत करणे
	\item थंड, अंधाऱ्या खोलीत ७ तास झोप = शरीराला पूर्ण विश्रांती
	\item जागरूकपूर्वक खाणे: काय खातो + \textbf{कधी खातो} --- दोन्ही महत्त्वाचे
	\item काळी कॉफी (साखर नाही, दूध नाही) योग्य प्रकारे घेतल्यास दीर्घायुष्यासाठी फायदेशीर
	\item कार्बोहायड्रेट्सच्या बाबतीत वेळ ही सर्व काही आहे --- बहुतेक आहारपद्धती हे सांगत नाहीत
	\end{itemize}

{\tiny (Ref: सी. शिवशंकरन --- दीर्घायुष्य)}

\end{frame}

%%%%%%%%%%%%%%%%%%%%%%%%%%%%%%%%%%%%%%%%%%%%%%%%%%%%%%%%%%%
\begin{frame}[fragile]\frametitle{व्यावहारिक बायोहॅक्स}

	\begin{itemize}
	\item सकाळी ३० सेकंद सूर्यप्रकाश घेणे
	\item आंघोळीचा शेवट थंड पाण्याने करणे
	\item कशासाठी कृतज्ञ आहात ते लिहून ठेवणे
	\item रक्त तपासणी करा (व्हिटॅमिन डी, मॅग्नेशियम, ओमेगा-३) --- अंदाज नको, डेटा पहा; तपासणीनंतरच पूरक आहार
	\item फक्त ५ किलो स्नायू वाढवा --- वृद्धत्व कसे येते यात मोठा फरक पडतो
	\item विशेषतः ४० च्या पुढे स्नायू = दीर्घायुष्याचा विमा
	\end{itemize}

{\tiny (Ref: सी. शिवशंकरन --- दीर्घायुष्य)}

\end{frame}

%%%%%%%%%%%%%%%%%%%%%%%%%%%%%%%%%%%%%%%%%%%%%%%%%%%%%%%%%%%
\begin{frame}[fragile]\frametitle{झोप + भारतीय स्वयंपाकघर}

	\begin{itemize}
	\item झोप ही महासत्ता आहे --- सर्व काही प्रभावित करते; फालतू वेळाची बर्बादी नव्हे
	\item झोपण्यापूर्वी: कीवी फळ किंवा ग्लायसिन घेणे; तोंडाला टेप लावणे; स्क्रीन लवकर बंद करणे
	\item आपल्या स्वयंपाकघरात उत्तरे आहेत:
		\begin{itemize}
		\item उडद डाळ (सालीसह), शेवगा, अंडी, रताळे, जांभळ्या रंगाच्या भाज्या
		\end{itemize}
	\item आपले आरोग्य संकट अटळ नाही --- ते सोडवता येण्यासारखे आहे
	\end{itemize}

{\tiny (Ref: सी. शिवशंकरन --- दीर्घायुष्य)}

\end{frame}
